In $\mathbb{Z}$ è possibile definire l'operazione di prodotto "$\cdot$" nel seguente modo

\[
    \forall a, b \in \mathbb{Z}, \quad a \cdot b := \sum_{k=1}^{b} a.
\]

Si può dimostrare che l'operazione di prodotto gode delle seguenti proprietà.

\begin{property}
    \textbf{Commutativa}
    \[\forall a, b \in \mathbb{Z}, \;a \cdot b = b \cdot a.\]
\end{property}
\begin{property}
    \textbf{Associativa}
    \[\forall a, b, c \in \mathbb{Z}, \;(a \cdot b) \cdot c = a \cdot (b \cdot c).\]
\end{property}
\begin{property}
    \textbf{Esistenza dell'elemento neutro}
    \[\forall a \in \mathbb{Z}, \; a \cdot 1 = 1 \cdot a = a.\]
\end{property}
\begin{property}
    \textbf{Chiusura di $\mathbb{Z}$ rispetto al prodotto}
    \[\forall a, b \in \mathbb{Z}, \; (a \cdot b) \in \mathbb{Z}.\]
\end{property}

\begin{observation}
    All'interno di $\mathbb{Z}$ non può essere definito l'elemento inverso rispetto al prodotto.
\end{observation}
\begin{observation}
    L'operazione di prodotto può essere applicata anche all'insieme $\mathbb{N}_0$, nei confronti della quale
    anche quest'ultimo risulta chiuso.
\end{observation}